\documentclass[11pt]{article}

\usepackage[T1]{fontenc}
\usepackage{amsmath,amssymb,amsthm}
\usepackage{booktabs}
\usepackage{array}
\usepackage{geometry}
\usepackage{graphicx}
\usepackage{natbib}
\usepackage{hyperref}
\usepackage{pgfplots}
\usepackage{xcolor}
\usepackage[expansion=false]{microtype}

\geometry{margin=1in}
\pgfplotsset{compat=1.18}
\hypersetup{colorlinks=true,linkcolor=black,citecolor=black,urlcolor=black}

\newtheorem{proposition}{Proposition}
\newtheorem{lemma}{Lemma}
\newtheorem{corollary}{Corollary}
\theoremstyle{remark}
\newtheorem{remark}{Remark}

\newcommand{\clip}{\operatorname{clip}}
\newcommand{\E}{\mathcal E}
\newcommand{\xh}{x_H}
\newcommand{\xu}{x_u}
\newcommand{\xb}{\bar x}

\title{Uniform-Price Equilibria and Welfare with Switching Costs:\\
A Reassessment of Gehrig, Shy, and Stenbacka (2012)}
\author{Ryota Matsuki\\
\small Independent Researcher, 790-0853 Matsuyama, Ehime, Japan\\
\small \texttt{ryota.matsuki@gmail.com}\\
\small ORCID: 0009-0005-2329-531X}
\date{}

\begin{document}

\maketitle

\begin{abstract}
We reassess the uniform-pricing benchmark in the post-referee accepted manuscript of \citet{GehrigShyStenbacka2012}. Retaining the complete clipped Hotelling demand reveals that the price vector reported in its Eq.~(12) is a pure-strategy Nash equilibrium only when inherited market dominance is sufficiently large. Writing \(s=\sigma/\tau\), the exact condition is
\[
x_0\ge \xh(s)
=\frac12-\frac{s}{3}
+\frac{\sqrt{3s(s+6)}}6.
\]
For \(0<s<1\), the reported price vector is the unique pure equilibrium on and above this boundary, while no pure-strategy price equilibrium exists below it. Mixed pricing is not characterized. The correction materially narrows the parameter set on which history-based and uniform pricing can be compared as pure equilibria. Reintegrating consumer utility over the actual allocation also resolves the weak-dominance consumer-surplus comparison: accepted Eq.~(15) is correct on its one-way-switch branch, but the transition to Eq.~(16) reverses its sign and the source integral is not valid on the no-switch branch. The reported sign reversal is not reproduced. On the corrected weak-dominance pure-equilibrium overlap, the smaller firm's profit is lower under history-based pricing, while the qualitative social-welfare ranking in favor of uniform pricing survives. Under strong dominance, the consumer-surplus comparison requires an explicit price-selection convention if below-cost poaching prices are admitted, whereas the welfare ranking is selection-invariant. Selected proof-critical algebraic and inequality statements are independently machine-checked in Lean 4/mathlib.
\end{abstract}

\begingroup\small
\noindent\textbf{Keywords:} history-based price discrimination; switching costs; Hotelling competition; pure-strategy equilibrium; consumer surplus; welfare.\\
\textbf{JEL:} D43, L13, L15.
\endgroup

\section{Introduction}
\label{sec:introduction}

Purchase-history data allow firms to distinguish inherited customers from consumers attached to a rival. When switching is costly, that information changes both the intensity of poaching and the value of an installed customer base. A large literature studies these forces in dynamic switching-cost competition and behavior-based price discrimination; see, among others, \citet{Klemperer1995}, \citet{Chen1997}, \citet{VillasBoas1999}, \citet{FudenbergTirole2000}, and \citet{ShafferZhang2000}.

Because customer-history information is the input that permits firms to distinguish retained customers from poaching targets, the welfare consequences of this form of data-enabled pricing depend on the correctness of the uniform-pricing benchmark. \citet{GehrigShyStenbacka2012} provide a particularly transparent welfare benchmark. Two Hotelling firms inherit asymmetric market shares, consumers pay a cost when changing supplier, and the firms are compared under uniform pricing and history-based price discrimination (HBP). In the post-referee accepted manuscript, the uniform-pricing analysis reports a unique Nash--Bertrand price vector and then uses that vector to compare market shares, profits, consumer surplus, and social welfare across pricing regimes.

This paper revisits that benchmark from the primitive consumer-choice problem. The key observation is that uniform demand is globally piecewise. Besides the smooth branch on which some customers of the inherited dominant firm switch to the smaller firm, there is a no-switch plateau. A local first-order solution on the switching branch is therefore not a Nash equilibrium unless each firm also prefers it to the relevant plateau boundary.

The resulting correction is exact. Let \(x_0>1/2\) be the dominant firm's inherited share, \(\tau\) the Hotelling transport parameter, and \(\sigma\) the switching cost. Writing \(s=\sigma/\tau\), define
\[
x_H(s)
=
\frac12-\frac{s}{3}
+\frac{\sqrt{3s(s+6)}}6.
\]
For \(0<s<1\), the uniform-price vector reported in the accepted manuscript is the unique pure-strategy equilibrium if and only if \(x_0\ge x_H(s)\). For \(1/2<x_0<x_H(s)\), no pure-strategy price equilibrium exists. The proof is global: it reconstructs the complete clipped demand, excludes capture and plateau equilibria, rules out the opposite switching branch, and compares the source candidate with the smaller firm's no-poaching kink. At equality, the smaller firm has two payoff-equal best responses, but the kink action does not form a second Nash profile.

Figure~\ref{fig:domains} visualizes the corrected equilibrium domain, while Table~\ref{tab:result-map} summarizes the downstream status of the accepted manuscript's numbered Results.

The correction has a first-order consequence for the welfare comparison. The accepted weak-HBP region and the pure-uniform equilibrium region overlap only if
\[
s<s_c
=
-3+\frac{6\sqrt{33}}{11}
\approx0.1334.
\]
Thus a substantial portion of the parameter space used for the source profile comparisons does not support the pure uniform equilibrium being compared.

A second correction is logically independent of equilibrium existence. At the reported uniform price vector, the realized allocation itself changes when the inherited share crosses
\[
x_u(s)=\frac12+\frac{s}{6}.
\]
Below this cutoff, no consumer switches. The one-way-switch consumer-surplus integral used in the accepted manuscript cannot be extrapolated through that ordering change. Direct primitive integration yields a piecewise consumer-surplus formula and shows that, throughout the accepted weak-dominance profile domain, HBP gives higher consumer surplus than the uniform source profile. Consequently, the weak-dominance consumer-surplus sign reversal in accepted Result~3 is not reproduced.

Combining the equilibrium and accounting corrections produces a compact downstream map. On the corrected weak pure-equilibrium overlap, the smaller firm earns less under HBP than under uniform pricing, so the high-switching-cost benefit region in accepted Result~4 is not an equilibrium comparison. The qualitative welfare ranking in accepted Result~5 survives: uniform pricing produces higher social welfare on the corrected overlap. Under strong dominance, the consumer-surplus result survives on the corrected pure domain under the source/nonnegative-margin HBP selection. If below-cost poaching prices are admitted, a zero-sales equilibrium price family makes consumer surplus and profit distribution selection-dependent. Social welfare is unchanged across that family, so the strong welfare ranking in accepted Result~7 survives independently of price selection.

The paper's contribution is deliberately model-specific. We do not claim a general nonexistence theorem for switching-cost Bertrand games, and we do not solve the mixed-price equilibrium below \(x_H\). Related dynamic, policy-choice, and imperfect-recognition models have different state variables or information structures \citep{JeongMaruyama2008,FabraGarcia2015,ColomboGrazianoPignataro2024}. The contribution here is to give the complete pure-price correspondence for the GSS uniform benchmark and to propagate that correction through the source welfare results without changing the underlying model.

The analysis is supported by three independent verification layers. Exact rational code evaluates demand directly from primitive clipped thresholds and searches global unilateral deviations; symbolic code reconstructs the surplus and welfare identities; and selected proof-critical algebra is machine-checked in Lean 4/mathlib. The Lean certificate is intentionally narrower than the economic theorem: the continuum demand partition and complete Nash branch enumeration remain analytic objects.

The remainder of the paper is organized as follows. Section~\ref{sec:model} states the source model and reconstructs exact uniform demand. Section~\ref{sec:uniform} derives the complete pure-strategy uniform-price correspondence and the corrected comparison domain. Section~\ref{sec:comparisons} recomputes market shares, consumer surplus, and profits. Section~\ref{sec:welfare} studies welfare and maps the implications to the accepted manuscript's numbered Results. Section~\ref{sec:scope} records strategy-domain, selection, and formal-verification boundaries. Sections~\ref{sec:literature}--\ref{sec:discussion} place the correction in the literature and discuss its interpretation.

\section{Model and source benchmark}
\label{sec:model}

We retain the static inherited-history Hotelling environment used in the accepted manuscript of \citet{GehrigShyStenbacka2012}. The purpose is not to introduce a new switching-cost model, but to resolve the global price game generated by those primitives.

\subsection{Consumers, histories, and switching costs}

Two firms, \(A\) and \(B\), are located at the endpoints \(0\) and \(1\) of a Hotelling line. A unit mass of consumers is uniformly distributed on \(z\in[0,1]\). Each consumer purchases one unit. Gross utility is \(\beta\), transport cost is linear with coefficient \(\tau>0\), and the two firms have a common marginal cost \(c\).

Consumers arrive with inherited supplier histories. Firm \(A\) inherits consumers on \([0,x_0]\), while firm \(B\) inherits consumers on \((x_0,1]\), where
\[
x_0\in(1/2,1].
\]
Thus \(A\) is the inherited dominant firm. A consumer who changes supplier pays a switching cost \(\sigma\), and throughout the reassessment we maintain the source restriction
\[
0\le \sigma<\tau.
\]

Under uniform prices \(p_A,p_B\), an \(A\)-history consumer located at \(z\le x_0\) obtains
\[
U_A^A(z)=\beta-p_A-\tau z,
\qquad
U_B^A(z)=\beta-p_B-\tau(1-z)-\sigma,
\]
whereas a \(B\)-history consumer located at \(z>x_0\) obtains
\[
U_A^B(z)=\beta-p_A-\tau z-\sigma,
\qquad
U_B^B(z)=\beta-p_B-\tau(1-z).
\]
Full coverage is assumed, so only the identity of the chosen firm matters for demand.

It is useful to normalize net margins and switching cost by the transport scale:
\[
a=\frac{p_A-c}{\tau},\qquad
b=\frac{p_B-c}{\tau},\qquad
s=\frac{\sigma}{\tau},\qquad
d=b-a.
\]
The normalization is one-to-one for fixed \(\tau>0\), and the equilibrium prices are recovered by multiplying normalized margins by \(\tau\) and adding \(c\).

\subsection{Primitive uniform demand}

The indifferent locations for the two inherited histories are
\[
t_A(d)=\frac{d+1+s}{2},
\qquad
t_B(d)=\frac{d+1-s}{2}.
\]
The first threshold determines which \(A\)-history consumers remain with \(A\); the second determines whether any \(B\)-history consumers switch to \(A\). Because each threshold is economically relevant only inside its own inherited segment, the exact share of firm \(A\) is
\begin{equation}
m_A(d)
=
\clip\!\left(t_A(d),0,x_0\right)
+
\clip\!\left(t_B(d)-x_0,0,1-x_0\right),
\label{eq:clipped-demand}
\end{equation}
where \(\clip(y,\ell,u)=\min\{\max\{y,\ell\},u\}\). Firm \(B\)'s share is \(m_B(d)=1-m_A(d)\).

For \(0<s<1\), define
\[
L=-1-s,\qquad
\alpha=2x_0-1-s,\qquad
\gamma=2x_0-1+s,\qquad
U=1+s.
\]
Equation~\eqref{eq:clipped-demand} is equivalent to the five-piece demand correspondence
\begin{equation}
m_A(d)=
\begin{cases}
0,&d\le L,\\[0.25em]
(d+1+s)/2,&L<d<\alpha,\\[0.25em]
x_0,&\alpha\le d\le\gamma,\\[0.25em]
(d+1-s)/2,&\gamma<d<U,\\[0.25em]
1,&d\ge U.
\end{cases}
\label{eq:five-piece-demand}
\end{equation}
The middle interval is the no-switch plateau. The lower switching branch has only \(A\)-history consumers moving to \(B\); the upper branch has only \(B\)-history consumers moving to \(A\). Simultaneous switching in both directions is impossible because \(t_A-t_B=s>0\).

Firm profits are
\[
\pi_A=\tau a\,m_A(d),\qquad
\pi_B=\tau b\,m_B(d).
\]
The common factor \(\tau\) is irrelevant for best responses.

\subsection{History-based prices used for comparison}

Under history-based price discrimination (HBP), let \(p_i\) denote firm \(i\)'s price to its inherited consumers and \(q_i\) its poaching price to the rival's inherited consumers. The accepted manuscript separates weak and strong inherited dominance at
\[
\bar x(s)=\frac{3-s}{4}.
\]

On the weak-dominance branch, \(1/2<x_0<\bar x(s)\), the source price vector is
\begin{align}
\frac{p_A^d-c}{\tau}&=\frac{2x_0+1+s}{3},
&
\frac{q_A^d-c}{\tau}&=\frac{3-4x_0-s}{3},
\nonumber\\
\frac{p_B^d-c}{\tau}&=\frac{3-2x_0+s}{3},
&
\frac{q_B^d-c}{\tau}&=\frac{4x_0-1-s}{3}.
\label{eq:weak-hbp-prices}
\end{align}
These prices and their induced allocation reproduce the accepted weak-HBP benchmark.

Under strong dominance, \(x_0\ge\bar x(s)\), the source member sets \(q_A^d=c\). For the nondegenerate case \(x_0<1\), allowing prices below marginal cost creates a payoff-equivalent zero-sales family discussed in Section~\ref{sec:scope}; the nonnegative-margin convention selects the source member. The endpoint \(x_0=1\) is different: the \(B\)-history market is empty, so \(q_A\) and \(p_B\) are unused prices and the HBP price vector is not unique even under nonnegative margins. Allocation, realized profits, consumer surplus, and welfare do not depend on those unused quotes.

\subsection{Source boundary}

All equation-level comparisons in this paper use a lawful post-referee accepted manuscript of \citet{GehrigShyStenbacka2012} as the controlling source. The exact Springer typeset Version of Record has not been compared equation by equation. We therefore refer to the accepted manuscript when identifying a formula or numbered Result whose wording matters.

\section{The uniform-price game}
\label{sec:uniform}

The accepted manuscript reports the normalized price pair
\[
a^*=1+\frac{s}{3},
\qquad
b^*=1-\frac{s}{3},
\]
and treats it as the unique Nash--Bertrand equilibrium in uniform prices. Those prices solve the first-order conditions on the lower switching branch of \eqref{eq:five-piece-demand}. Global optimality, however, also requires comparison with the no-switch kink.

Define
\[
x_u(s)=\frac12+\frac{s}{6}
\]
and
\begin{equation}
x_H(s)
=
\frac12-\frac{s}{3}
+\frac{\sqrt{3s(s+6)}}{6}.
\label{eq:xH}
\end{equation}
For every \(0<s<1\),
\[
\frac12<x_u(s)<x_H(s)<1.
\]

\begin{proposition}[Complete pure uniform-price correspondence]
\label{prop:uniform}
Let \(x_0\in(1/2,1]\), \(\tau>0\), \(c\in\mathbb R\), and \(0\le\sigma<\tau\). For \(s=0\), the unique pure uniform-price equilibrium has normalized margins \((a,b)=(1,1)\). For \(0<s<1\),
\[
\E^u(x_0,s)
=
\begin{cases}
\left\{\left(1+s/3,\,1-s/3\right)\right\},
&x_0\ge x_H(s),\\[0.4em]
\varnothing,
&1/2<x_0<x_H(s).
\end{cases}
\]
Equivalently, whenever the pure equilibrium exists,
\[
p_A^u=c+\tau+\frac{\sigma}{3},
\qquad
p_B^u=c+\tau-\frac{\sigma}{3}.
\]
At \(x_0=x_H(s)\), firm \(B\) has an additional payoff-equal kink best response against \(A\)'s equilibrium price, but that kink action does not generate a second Nash profile.
\end{proposition}

\paragraph{Mechanism.}
On the lower switching branch, the two smooth first-order conditions have the unique intersection above. Firm \(A\)'s source action is globally optimal once the candidate lies on that branch. Firm \(B\), however, can instead move to the right edge of the no-switch plateau and raise its price while retaining all inherited \(B\)-history consumers. Against \(a^*\), this no-poaching kink price is
\[
b_k=a^*+\gamma=2x_0+\frac{4s}{3}.
\]
Let
\[
x_L(s)=
\frac12-\frac{s}{3}
-\frac{\sqrt{3s(s+6)}}{6}.
\]
The exact payoff comparison factors as
\begin{equation}
\frac{\pi_B(a^*,b^*)-\pi_B(a^*,b_k)}{\tau}
=
2\bigl(x_0-x_L(s)\bigr)\bigl(x_0-x_H(s)\bigr).
\label{eq:kink-factor}
\end{equation}
Since \(x_L(s)<1/2<x_0\), firm \(B\)'s source price is globally optimal if and only if \(x_0\ge x_H(s)\). The complete exclusion of capture, plateau, kink, and upper-switching alternatives is given in Appendix~\ref{app:uniform-proof}.

\begin{remark}[Pure rather than general Nash nonexistence]
For \(1/2<x_0<x_H(s)\), Proposition~\ref{prop:uniform} establishes the nonexistence of a \emph{pure-strategy} price equilibrium. We do not characterize mixed pricing and make no statement that Nash equilibrium itself fails to exist.
\end{remark}

\subsection{An exact counterexample}

Take
\[
\tau=1,\qquad
\sigma=\frac12,\qquad
x_0=\frac35,\qquad
c=0.
\]
The accepted source candidate is
\[
p_A=\frac76,\qquad p_B=\frac56.
\]
Firm \(B\)'s profit is \(25/72\). Moving to the no-poaching kink
\[
p_B'=\frac{28}{15}
\]
raises its profit to \(56/75\), so the exact gain is
\[
\frac{56}{75}-\frac{25}{72}
=
\frac{719}{1800}>0.
\]
This point lies inside the accepted weak-dominance domain, so the issue is a global best-response failure rather than an inadmissible parameter choice.

\begin{table}[t]
\centering
\caption{Exact global-deviation counterexample at \(\tau=1\), \(\sigma=1/2\), \(x_0=3/5\), and \(c=0\).}
\label{tab:counterexample}
\begin{tabular}{lccc}
\toprule
 & \(p_B\) & \(m_B\) & \(\pi_B\)\\
\midrule
Accepted Eq.~(12) candidate & \(5/6\) & \(5/12\) & \(25/72\)\\
No-poaching kink deviation & \(28/15\) & \(2/5\) & \(56/75\)\\
\bottomrule
\end{tabular}

\smallskip
\footnotesize
The deviation gain is \(719/1800\). Quantities are calculated from the primitive clipped demand.
\end{table}

\subsection{The corrected comparison domain}

The accepted weak-HBP domain is \(1/2<x_0<\bar x(s)\), where \(\bar x(s)=(3-s)/4\). Comparing it with Proposition~\ref{prop:uniform} yields an overlap only when
\begin{equation}
0\le s<s_c,
\qquad
s_c=-3+\frac{6\sqrt{33}}{11}
\approx0.1333978.
\label{eq:sc}
\end{equation}
For \(0<s<s_c\), two pure equilibria can be compared only on
\[
x_H(s)\le x_0<\bar x(s).
\]
At \(s=s_c\), the two boundaries coincide and the weak-HBP inequality is strict, so the overlap is empty.

Figure~\ref{fig:domains} plots the three relevant boundaries. The figure is generated from the verified Stage-9 data file rather than from hand-entered curves.

\begin{figure}[t]
\centering
\begin{tikzpicture}
\begin{axis}[
  width=0.88\textwidth,
  height=0.50\textwidth,
  xlabel={normalized switching cost \(s=\sigma/\tau\)},
  ylabel={inherited share \(x_0\)},
  xmin=0,xmax=1,
  ymin=0.5,ymax=1,
  legend style={draw=none,at={(0.5,-0.20)},anchor=north,legend columns=1},
  grid=major
]
\addplot+[name path=xH,thick,mark=none] table[x=s,y=x_H,col sep=comma]{generated/parameter_domains.csv};
\addlegendentry{\(x_H(s)\): pure-uniform threshold}
\addplot+[name path=xbar,dashed,thick,mark=none] table[x=s,y=x_bar,col sep=comma]{generated/parameter_domains.csv};
\addlegendentry{\(\bar x(s)\): weak/strong HBP boundary}
\addplot+[dotted,thick,mark=none] table[x=s,y=x_u,col sep=comma]{generated/parameter_domains.csv};
\addlegendentry{\(x_u(s)\): uniform switching cutoff}
\addplot[fill opacity=0.12,draw=none,forget plot]
  fill between[of=xbar and xH,soft clip={domain=0:0.1333978072}];
\addplot+[dashed,mark=none,forget plot] coordinates {(0.1333978072,0.5) (0.1333978072,1.0)};
\addplot+[only marks,mark=*,mark size=1.8pt,forget plot] coordinates {(0.1333978072,0.7166505482)};
\node[anchor=west] at (axis cs:0.145,0.94) {\(s_c\)};
\end{axis}
\end{tikzpicture}
\caption{Corrected equilibrium and source-branch boundaries. The shaded region is the weak-HBP/pure-uniform equilibrium overlap \(x_H(s)\le x_0<\bar x(s)\), which exists only for \(s<s_c\); the vertical dashed line marks \(s_c\).}
\label{fig:domains}
\end{figure}

\section{Consumer surplus and profit comparisons}
\label{sec:comparisons}

The equilibrium correction changes the domain on which the two pricing regimes can be compared, but it does not by itself determine the sign of consumer surplus or profits. Those objects must be recomputed from the realized allocation.

\subsection{Market shares}

At the uniform candidate, the relevant \(A\)-history cutoff is
\[
x_u=\frac12+\frac{s}{6}.
\]
Whenever the uniform profile is a pure equilibrium, Proposition~\ref{prop:uniform} implies \(x_0\ge x_H(s)>x_u(s)\). Thus only \(A\)-history consumers in \((x_u,x_0]\) switch to \(B\), and firm \(A\)'s equilibrium share is \(x_u\).

On the accepted weak-HBP branch, firm \(A\)'s share is \((2-x_0)/3\). Hence, whenever the two pure equilibria coexist,
\begin{equation}
x_u-\frac{2-x_0}{3}
=
\frac{2x_0+s-1}{6}>0.
\label{eq:share-gap}
\end{equation}
The qualitative market-share ranking in accepted Result~2 therefore survives, although the weak-branch displayed difference in the accepted manuscript does not match its own equilibrium shares.

\subsection{Weak-dominance consumer surplus}

The uniform profile has two possible allocation branches even before one asks whether it is an equilibrium. If \(x_0\ge x_u\), some \(A\)-history consumers switch to \(B\). If \(x_0<x_u\), no consumer switches and the inherited allocation remains unchanged. Integrating realized utility over the clipped allocation gives the following profile comparison.

\begin{proposition}[Branch-correct weak-HBP consumer surplus]
\label{prop:weak-cs}
For the accepted weak-HBP price profile and the accepted uniform-price profile,
\[
CS^d-CS^u=
\begin{cases}
\displaystyle
\frac{-\tau^2(52x_0^2-52x_0-1)
+2\sigma\tau(18x_0-17)+\sigma^2}{36\tau},
&x_0\ge x_u,\\[1.2em]
\displaystyle
\frac{\sigma^2+12\sigma\tau x_0-14\sigma\tau
-8\tau^2x_0^2+8\tau^2x_0+5\tau^2}{18\tau},
&x_0<x_u.
\end{cases}
\]
The two expressions agree at \(x_0=x_u\), and \(CS^d-CS^u>0\) throughout the accepted weak-dominance profile domain \(1/2<x_0<\bar x(s)\).
\end{proposition}

The first branch differs algebraically from accepted Eq.~(15). More importantly, the accepted one-way-switch integral cannot be extended into \(x_0<x_u\), because its switching interval is then reversed. The exact point \((s,x_0)=(0.99,0.501)\), for example, lies on the no-switch branch and gives
\[
CS^d-CS^u=\frac{17993}{4500000}>0.
\]
At a valid switching-branch point inside the corrected pure-equilibrium overlap, \((s,x_0)=(0.1,0.7)\), direct integration yields \(221/720\), while the accepted Eq.~(15) expression evaluates to \(1279/3600\). Both are positive, but they are not equal.

\begin{corollary}[Weak-HBP equilibrium consumer surplus]
\label{cor:weak-cs}
Whenever weak HBP and uniform pricing both admit the pure equilibria being compared,
\[
CS^d>CS^u.
\]
Thus the consumer-surplus sign reversal stated in accepted Result~3 is not reproduced on the corrected pure-equilibrium domain.
\end{corollary}

\subsection{The smaller firm's profit under weak dominance}

Using the accepted weak-HBP prices and the corrected uniform equilibrium prices, the normalized profit difference for firm \(B\) is
\begin{equation}
\frac{\pi_B^d-\pi_B^u}{\tau}
=
\frac{s^2-12sx_0+14s+20x_0^2-20x_0+1}{18}.
\label{eq:B-profit-gap}
\end{equation}

\begin{proposition}[Smaller-firm profit on the pure overlap]
\label{prop:B-profit}
For every \(0\le s<s_c\) and every \(x_H(s)\le x_0<\bar x(s)\),
\[
\pi_B^d<\pi_B^u.
\]
\end{proposition}

The numerator of \eqref{eq:B-profit-gap} is convex in \(x_0\). On the larger interval \([1/2,\bar x]\), its endpoint values are
\[
s^2+8s-4<0
\quad\text{and}\quad
\frac{(1+s)(21s-11)}4<0
\]
for \(0\le s<s_c\). Convexity then implies negativity throughout the interval. Consequently, the high-switching-cost region associated with accepted Result~4 can still be evaluated against the Eq.~(12) counterfactual profile, but it is not a comparison of two pure equilibria.

\subsection{What the equilibrium correction does}

The distinction between profile and equilibrium comparisons is central. Proposition~\ref{prop:weak-cs} is a statement about two specified price profiles on the entire source weak-dominance domain. Corollary~\ref{cor:weak-cs} and Proposition~\ref{prop:B-profit} are equilibrium statements only on
\[
0<s<s_c,
\qquad
x_H(s)\le x_0<\bar x(s).
\]
No consumer-surplus or profit value is attributed to an uncharacterized mixed uniform-pricing equilibrium below \(x_H\).

\section{Social welfare and the accepted Results}
\label{sec:welfare}

Because the market is fully covered and each consumer buys one unit, prices are transfers in aggregate welfare. Let \(y(z)\in\{0,1\}\) denote the chosen firm's location and let \(N_{\rm sw}\) denote the mass of consumers who switch away from their inherited supplier. Then
\begin{equation}
W=CS+\pi_A+\pi_B
=
\beta-c
-\tau\int_0^1 |z-y(z)|\,dz
-\sigma N_{\rm sw}.
\label{eq:welfare-resource}
\end{equation}
Thus the welfare comparison can be checked directly from real transport and switching costs, independently of price-transfer accounting.

\subsection{Weak dominance}

On the only uniform allocation branch compatible with a weak-domain pure equilibrium,
\begin{equation}
W^d-W^u
=
\frac{\tau}{36}
\left[
28x_0^2-28x_0+5
+2s(18x_0-13)+5s^2
\right].
\label{eq:weak-welfare}
\end{equation}

\begin{proposition}[Weak-dominance welfare]
\label{prop:weak-welfare}
On the corrected weak-HBP/pure-uniform overlap,
\[
W^d<W^u.
\]
\end{proposition}

Expression~\eqref{eq:weak-welfare} is increasing in \(x_0>1/2\). At the weak-dominance upper boundary \(x_0=\bar x=(3-s)/4\),
\[
W^d-W^u
=
-\frac{\tau(1+s)(1+9s)}{144}<0.
\]
The equilibrium overlap lies weakly below that boundary, so the gap is strictly negative throughout. The accepted manuscript's endpoint substitution following its weak-welfare formula is algebraically incorrect, but the qualitative welfare conclusion in Result~5 survives after imposing the corrected equilibrium domain.

\subsection{Strong dominance}

For the source strong-HBP member \(q_A^d=c\), and on the corrected uniform pure-equilibrium domain, direct integration gives
\begin{equation}
CS^d-CS^u
=
\frac{\tau(1-x_0)(16-7x_0-13s)}9.
\label{eq:strong-cs}
\end{equation}
Therefore the accepted strong-dominance consumer-surplus threshold survives under the source/nonnegative-margin HBP selection, conditional on the corrected uniform pure-existence condition.

Social welfare is cleaner. Across the entire strong-HBP zero-sales price family described in Section~\ref{sec:scope}, the allocation is unchanged. On the corrected uniform pure-equilibrium domain,
\begin{equation}
W^d-W^u
=
-\frac{\tau(1-x_0)(1-x_0+2s)}9.
\label{eq:strong-welfare}
\end{equation}

\begin{proposition}[Strong-dominance welfare]
\label{prop:strong-welfare}
For \(x_0\ge\bar x(s)\) such that the uniform pure equilibrium exists,
\[
W^d\le W^u,
\]
with strict inequality for \(x_0<1\) and equality only at \(x_0=1\). This welfare ranking is invariant to the strong-HBP below-cost zero-sales price selection.
\end{proposition}

\subsection{Impact on the accepted manuscript}

Table~\ref{tab:result-map} summarizes the corrected status of the seven numbered Results in the accepted manuscript.

\begin{table}[t]
\centering
\caption{Impact on numbered Results in the accepted manuscript}
\label{tab:result-map}
\small
\begin{tabular}{p{0.10\textwidth}p{0.32\textwidth}p{0.50\textwidth}}
\toprule
Result & Subject & Reassessment\\
\midrule
1 & Switching costs above the maintained \(\sigma<\tau\) domain & Outside the maintained reassessment domain; no correction claim. \\
2 & Dominant firm's market share & Qualitative sign survives; the weak-branch displayed difference is inconsistent with the source shares. Equilibrium wording requires \(x_0\ge x_H\). \\
3 & Weak consumer surplus & Does not survive as stated. Branch-correct primitive integration gives \(CS^d>CS^u\) throughout the weak source profile domain. \\
4 & Smaller-firm profit & The source high-switching-cost benefit region does not overlap the corrected weak pure-equilibrium domain; on that overlap \(\pi_B^d<\pi_B^u\). \\
5 & Weak social welfare & Qualitative ranking in favor of uniform pricing survives on the corrected pure overlap; the source endpoint substitution is algebraically incorrect. \\
6 & Strong consumer surplus & Survives on the corrected pure domain under the source/nonnegative-margin strong-HBP selection; unrestricted below-cost pricing introduces selection dependence. \\
7 & Strong social welfare & Survives on the corrected pure domain and is selection-invariant; equality occurs at \(x_0=1\). \\
\bottomrule
\end{tabular}
\end{table}

\section{Strategy domains, robustness, and formal verification}
\label{sec:scope}

\subsection{Uniform-price strategy domain}

Proposition~\ref{prop:uniform} is unchanged whether uniform prices are allowed on all of \(\mathbb R\) or restricted to \(p_i\ge c\). A negative-margin price with positive demand yields negative profit and can be replaced by \(c\), while a zero-margin price with positive demand can be raised slightly without losing all demand. Zero-demand capture profiles are excluded by profitable entry deviations. This reduction is specific to the uniform-price game and should not be silently transferred to the history-based game.

\subsection{Strong-HBP price selection}

Under strong dominance with \(x_0<1\), the accepted source member has \(q_A^d=c\). If below-cost poaching prices are admissible, however, there is a bounded family of payoff-equivalent zero-sales actions. Let
\[
u=\frac{q_A^d-c}{\tau}.
\]
Then
\begin{equation}
3-s-4x_0\le u\le0,
\qquad
\frac{p_B^d-c}{\tau}=u+2x_0-1+s.
\label{eq:strong-family}
\end{equation}
Across this family, firm \(A\)'s poaching price still attracts no \(B\)-history consumers. The allocation and real resource costs are unchanged, but the price paid by retained \(B\)-history consumers changes. Consumer surplus and the distribution of firm profits therefore depend on \(u\), while total welfare does not.

For \(\bar x(s)\le x_0<1\), the nonnegative-margin restriction \(q_A^d\ge c\) collapses the interval in \eqref{eq:strong-family} to the source member \(u=0\). Equation~\eqref{eq:strong-cs} and the consumer-surplus interpretation of accepted Result~6 use that selection. Proposition~\ref{prop:strong-welfare} does not require it.

At \(x_0=1\), the \(B\)-history segment has zero mass. The quotes \(q_A\) and \(p_B\), which apply only to that segment, are payoff-irrelevant and may be varied without changing any realized choice. Thus neither unrestricted prices nor nonnegative margins deliver a unique four-price vector at this endpoint. The active-market prices \(p_A\) and \(q_B\), the allocation, both firms' realized profits, consumer surplus, and social welfare are invariant. In particular, the strong-HBP/uniform share and consumer-surplus differences are zero at \(x_0=1\).

\subsection{Normalization and portability}

The exact threshold \(x_H\) is invariant to a common translation in marginal cost \(c\) and to the transport-scale normalization. In dimensional prices, the equilibrium vector is always
\[
p_A^u=c+\tau+\sigma/3,
\qquad
p_B^u=c+\tau-\sigma/3.
\]
The direct clipped-demand representation and the price-difference representation also yield the same correspondence, and the allocation convention for a measure-zero set of indifferent consumers is immaterial.

These invariances do not turn Proposition~\ref{prop:uniform} into a theorem about arbitrary switching-cost Bertrand games. The proof uses the endpoint Hotelling geometry, the inherited split at \(x_0\), linear transport costs, full coverage, and the resulting five-piece demand. The contribution is deliberately model-specific.

\subsection{Formal verification boundary}

Selected proof-critical algebra has been independently checked in Lean 4.19.0 with mathlib pinned to commit
\texttt{c44e0c8ee63ca166450922a373c7409c5d26b00b}.
Formal coverage is classified as \textbf{PROOF-CRITICAL CORE}. The machine-checked targets include the payoff factorization in \eqref{eq:kink-factor}, \(x_u<x_H<1\), the exact \(719/1800\) counterexample gain, the equality-kink deviation gain, representative consumer-surplus identities, the weak-welfare endpoint identity, and the polynomial characterization of \(s_c\).

The Lean development does not encode the continuum consumer model, the complete clipped demand correspondence, the Nash definition, or the global enumeration of equilibrium branches. Those objects are established analytically and by an independent direct evaluator. Formal verification therefore covers the proof-critical algebraic core, not the complete economic model.

\subsection{Interpretation}

The model isolates how inherited customer histories and switching costs create a discrete global pricing option: a firm can stop poaching and move to a plateau boundary while retaining its inherited consumers. The analysis is theoretical. We make no empirical causal claim, calibration claim, or policy claim beyond the comparative welfare statements inside the maintained source model.

\section{Related literature}
\label{sec:literature}

This reassessment sits at the intersection of switching-cost competition and behavior- or history-based pricing. Switching costs make inherited market shares strategically valuable because firms trade off harvesting attached consumers against competing for rivals' consumers; \citet{Klemperer1995} surveys this broad mechanism. In dynamic models, customer acquisition affects future states and prices. Examples include \citet{Chen1997}, \citet{VillasBoas1999}, and, in a continuous-time setting with asymmetric market shares, \citet{FabraGarcia2015}. Our result is not a dynamic switching-cost theorem: inherited shares are fixed primitives of the one-period pricing game considered by \citet{GehrigShyStenbacka2012}.

A second strand studies customer recognition and targeted poaching. \citet{FudenbergTirole2000} analyze customer poaching and brand switching under alternative contract structures, while \citet{ShafferZhang2000} study preference-based discrimination with switching costs. \citet{AcquistiVarian2005} examine conditioning prices on purchase histories, including competitive environments. These papers establish that customer information can intensify competition and alter switching, but they do not supply the global best-response condition in the particular inherited-history Hotelling benchmark studied here.

A particularly close recent neighborhood combines asymmetric firms, switching costs, and behavior-based pricing. \citet{Umezawa2022} studies a two-period horizontally and vertically differentiated duopoly. An author-issued corrigendum, \citet{UmezawaCorrigendum2023}, is methodologically relevant here: it corrects the uniform-pricing benchmark because a price ordering had been used to infer a switching/allocation ordering that need not follow. The corrected analysis classifies regimes by the actual allocation cutoffs instead. This is a direct precedent for treating regime consistency as a mathematical obligation rather than as a consequence of a candidate price ordering.

The Umezawa model is nevertheless not identified here as a parent theorem for the present correction. Its timing is two-period, horizontal and vertical differentiation jointly determine demand, and customer histories are generated dynamically. The present paper instead takes the GSS inherited split as a primitive and solves a one-period uniform-price game. Likewise, \citet{Shrivastav2023} changes customer recognition through imperfect information and misinformation; \citet{ColomboGrazianoPignataro2024} studies imperfect history-based recognition with asymmetric market shares; and \citet{UmezawaYamakawa2025} introduces multiple consumer types and a different switching-cost incidence. We did not establish a formal nesting or non-nesting theorem for these models. The narrower claim is that our search located no prior derivation of the GSS-specific threshold \(x_H(s)\) or its complete pure-price correspondence.

The closest policy-choice comparison is \citet{JeongMaruyama2008}, where firms choose between uniform and discriminatory pricing under exogenous switching costs. Their equilibrium object includes the pricing-policy choice itself, whereas our object is the complete pure-strategy correspondence of the fixed uniform-price subgame in the GSS benchmark.

Accordingly, the contribution is narrower than a new general theory of behavior-based discrimination. It is a mathematical reassessment of a specific benchmark that is frequently interpretable within this literature: the accepted uniform-price solution is a local switching-branch candidate, and its global pure-equilibrium validity requires comparing it with a no-poaching kink. The resulting threshold changes which downstream welfare comparisons are equilibrium comparisons at all.

\section{Discussion}
\label{sec:discussion}

\subsection{Why the local first-order conditions are insufficient}

The accepted uniform-price vector is not generated by an algebraic mistake in the smooth first-order conditions. It is the correct intersection of those conditions on the lower switching branch. The problem is global. Demand contains a flat no-switch region. For the smaller firm, moving from the smooth poaching optimum to the right-hand plateau boundary can sharply increase price without sacrificing its inherited customers. The profitability of that finite deviation is exactly what \(x_H\) measures.

This distinction matters methodologically. In piecewise price games, verifying first-order and second-order conditions on one regime does not establish Nash equilibrium unless all feasible regime changes and kinks are checked. Here, the omitted comparison is economically meaningful: it is the choice between continuing to poach the dominant firm's customers and ceasing poaching while harvesting one's own inherited base.

\subsection{Why the welfare correction is not simply the equilibrium correction}

Two logically separate issues arise downstream. First, the uniform-price vector is an equilibrium only on the corrected domain. Second, even as a counterfactual price profile, its allocation changes at \(x_0=x_u\). Extending a one-way-switch consumer-surplus integral below that cutoff reverses the order of an integration interval and no longer represents realized consumer choices.

Consequently, one cannot repair the accepted analysis by adding the condition \(x_0\ge x_H\) alone. Consumer surplus must also be integrated over the actual clipped allocation. This is why the weak consumer-surplus reversal disappears even at the profile level, while the welfare ranking favoring uniform pricing survives on the equilibrium overlap.

\subsection{Economic implications inside the model}

The corrected overlap between weak HBP and a pure uniform-price equilibrium is small: normalized switching cost must satisfy \(s<s_c\approx0.1334\). In that region, HBP raises consumer surplus but lowers the smaller firm's profit and lowers total welfare relative to uniform pricing. The welfare effect reflects real switching and transport costs, not price transfers.

Under strong inherited dominance, the source consumer-surplus comparison is more sensitive to price-domain conventions because below-cost poaching prices can support the same zero-sales allocation. Social welfare is more robust: the entire zero-sales price family induces the same allocation, so the strong welfare ranking is selection-invariant.

\subsection{Limitations}

The analysis does not solve the mixed-strategy uniform-price game below \(x_H\). It therefore does not compare equilibrium consumer surplus or welfare in that region. Nor does it claim robustness to alternative spatial demand, outside options, heterogeneous switching costs, or imperfect recognition. Those are distinct models rather than omitted cases of the theorem.

A final source limitation is textual rather than mathematical. Equation-level statements are verified against a post-referee accepted manuscript. Until the exact typeset Version of Record is compared directly, this paper does not claim that every equation number or algebraic display in the final publisher layout is identical to that accepted version.

\section{Conclusion}
\label{sec:conclusion}

Reconstructing the uniform-price demand from primitive consumer choices changes the equilibrium domain in the inherited-history Hotelling benchmark of \citet{GehrigShyStenbacka2012}. The accepted Eq.~(12) price vector is the correct smooth switching-branch candidate, but it is a pure-strategy Nash equilibrium only when
\[
x_0\ge
\frac12-\frac{s}{3}
+\frac{\sqrt{3s(s+6)}}6.
\]
Below this boundary no pure-strategy price equilibrium exists; mixed pricing remains outside the scope of the reassessment.

The correction materially narrows the set of parameters over which weak HBP and uniform pricing can be compared as pure equilibria. Direct primitive integration also yields a separate consumer-surplus correction: accepted Eq.~(15) is correct on its switching branch, but Eq.~(16) reverses its sign and the no-switch region requires a different realized-allocation expression. The weak-dominance sign reversal is therefore not reproduced. On the corrected weak pure overlap, HBP raises consumer surplus, lowers the smaller firm's profit, and lowers social welfare relative to uniform pricing. Under strong dominance, consumer-surplus conclusions require an explicit source/nonnegative-margin price selection if below-cost poaching prices are allowed, whereas the welfare comparison is invariant to that selection.

The broader lesson is specific but useful. In price games with history-dependent kinks, an interior candidate should not be labeled an equilibrium until finite deviations across allocation regimes have been ruled out. Here that global check produces an exact threshold, changes the feasible welfare-comparison domain, and separates source conclusions that survive from those that do not.


\appendix
\section{Global proof of Proposition~\ref{prop:uniform}}
\label{app:uniform-proof}

This appendix records the complete pure-strategy branch argument. Normalize \(\tau=1\) and work with net margins \(a=(p_A-c)/\tau\) and \(b=(p_B-c)/\tau\). The unrestricted-real strategy-space proof and the nonnegative-margin proof are handled together.

\subsection*{Preliminary margin--demand elimination}

At any pure Nash equilibrium, each firm must have strictly positive demand and a strictly positive margin.

First, if a firm serves positive demand at a negative margin, choosing the zero margin yields payoff zero and strictly improves on its negative profit. If a firm serves positive demand at zero margin, continuity of the clipped demand implies that a sufficiently small positive increase in its margin preserves positive demand and yields strictly positive profit.

Second, suppose firm \(A\) has zero demand, so firm \(B\) serves the whole market. If \(B\)'s margin is negative, \(B\) can move to zero margin and improve to zero profit. If its margin is zero, a sufficiently small price increase preserves positive demand and yields positive profit. If its margin is strictly positive, \(A\) can choose the same positive margin as \(B\), which sets \(d=b-a=0\) and gives \(A\) strictly positive demand and profit. Hence no zero-demand profile can be Nash. The argument is symmetric for \(B\).

This joint margin--demand argument excludes the capture tails without assuming that a zero-demand firm can always enter at a positive margin against an arbitrary rival price.

\subsection*{The case \(s=0\)}

When \(s=0\), inherited histories do not affect demand:
\[
m_A(d)=\clip\left(\frac{d+1}{2},0,1\right).
\]
After the capture tails are excluded by the preliminary argument, each firm's payoff is a concave quadratic on the interior demand region. The two first-order conditions intersect uniquely at
\[
a=b=1.
\]
Thus the unique pure equilibrium for \(s=0\) is \((1,1)\). The remainder of the proof assumes \(0<s<1\).

\subsection*{The five pieces for \(0<s<1\)}

For fixed rival price, unilateral payoffs as functions of \(d=b-a\) are
\[
\begin{array}{c|c|c}
\text{region} & \pi_A(d;b) & \pi_B(d;a)\\
\hline
d\le L & 0 & a+d\\
L<d<\alpha &
(b-d)(d-L)/2 &
(a+d)(1-s-d)/2\\
\alpha\le d\le\gamma &
x_0(b-d) &
(1-x_0)(a+d)\\
\gamma<d<U &
(b-d)(d+1-s)/2 &
(a+d)(1+s-d)/2\\
d\ge U &
b-d & 0
\end{array}
\]
where \(L=-1-s\), \(\alpha=2x_0-1-s\), \(\gamma=2x_0-1+s\), and \(U=1+s\). Each switching-branch payoff is a concave quadratic and each plateau payoff is affine. A global best response therefore lies at a feasible quadratic vertex or a regime endpoint.

\paragraph{No-switch plateau and kinks.}
For \(x_0<1\), both firms have positive inherited shares on the plateau. Firm \(A\)'s plateau payoff \(x_0(b-d)\) strictly decreases in \(d\), while firm \(B\)'s payoff \((1-x_0)(a+d)\) strictly increases in \(d\). Hence no plateau interior can be Nash.

At \(d=\alpha\) with \(x_0<1\), firm \(B\) can move to \(d=\gamma\), retain share \(1-x_0\), and increase its margin by \(2s\), gaining \(2s(1-x_0)>0\). If \(x_0=1\), firm \(B\) has zero demand at \(d=\alpha\), so the preliminary zero-demand argument excludes the profile instead.

At \(d=\gamma\), firm \(A\) can move to \(d=\alpha\), retain share \(x_0\), and increase its margin by \(2s\), gaining \(2sx_0>0\). Thus neither kink is a Nash profile, including the endpoint \(x_0=1\).

\paragraph{Lower switching branch.}
For \(L<d<\alpha\), first-order conditions are
\[
a=d+1+s,
\qquad
b=1-s-d.
\]
Together with \(d=b-a\), these yield the unique smooth candidate
\[
d^*=-\frac{2s}{3},
\qquad
a^*=1+\frac{s}{3},
\qquad
b^*=1-\frac{s}{3}.
\]
The candidate lies strictly inside the lower switching branch exactly when
\[
x_0>x_u(s)=\frac12+\frac{s}{6}.
\]
At \(x_0=x_u\) it lies at the lower kink \(d=\alpha\), already excluded.

Against \(b^*\), firm \(A\)'s lower-switching payoff has its unique vertex at \(d^*\). Its best plateau payoff is at \(d=\alpha\), and the exact gap is
\[
\frac{\pi_A(a^*,b^*)-\pi_A(\alpha;b^*)}{\tau}
=
2(x_0-x_u)^2>0
\]
for \(x_0>x_u\). On the upper switching branch,
\[
\frac{\partial \pi_A}{\partial d}
=
\frac{b^*-2d-1+s}{2}.
\]
At its left endpoint \(d=\gamma\), this derivative equals
\[
1-2x_0-\frac{2s}{3}<0,
\]
and concavity makes it still smaller for larger \(d\). Hence the upper switching branch is maximized at \(\gamma\), whose plateau payoff is below the payoff at \(\alpha\). The capture tail gives zero. Thus \(a^*\) is firm \(A\)'s unique global best response to \(b^*\) whenever \(x_0>x_u\).

Against \(a^*\), firm \(B\)'s best plateau action is the no-poaching kink \(d=\gamma\), with
\[
b_k=a^*+\gamma
=
2x_0+\frac{4s}{3}.
\]
On the upper switching branch,
\[
\frac{\partial \pi_B}{\partial d}
=
\frac{1+s-a^*-2d}{2},
\]
which at \(d=\gamma\) again equals
\[
1-2x_0-\frac{2s}{3}<0.
\]
Thus the upper branch decreases from the kink and the right capture tail gives zero. On the left capture tail, \(B\)'s margin satisfies
\[
b=a^*+d\le a^*+L=-\frac{2s}{3}<0,
\]
so it cannot dominate the positive-profit source action. The relevant global comparison is therefore exactly
\[
\frac{\pi_B(a^*,b^*)-\pi_B(a^*,b_k)}{\tau}
=
2\bigl(x_0-x_L(s)\bigr)\bigl(x_0-x_H(s)\bigr).
\]
Since \(x_L(s)<1/2<x_0\), firm \(B\)'s source action is globally optimal if and only if \(x_0\ge x_H(s)\).

At equality \(x_0=x_H(s)\), \(B\) is indifferent between \(b^*\) and \(b_k\). But the kink profile \((a^*,b_k)\) is not Nash: firm \(A\) can move the price difference from \(\gamma\) to \(\alpha\), keep share \(x_0\), raise its margin by \(2s\), and gain \(2sx_0>0\). Thus best-response multiplicity for \(B\) does not generate a second Nash profile.

\paragraph{Upper switching branch.}
If both firms' first-order conditions were satisfied on \(\gamma<d<U\), they would imply
\[
d=\frac{2s}{3},
\qquad
a=1-\frac{s}{3},
\qquad
b=1+\frac{s}{3}.
\]
Feasibility would require
\[
\frac{2s}{3}>\gamma=2x_0-1+s
\quad\Longleftrightarrow\quad
x_0<\frac12-\frac{s}{6},
\]
which is impossible because \(x_0>1/2\). Hence no upper-switching interior equilibrium exists.

\subsection*{Conclusion}

For \(0<s<1\), the plateau, both kinks, both capture tails, and the upper-switching interior have been excluded. The lower-switching smooth candidate is globally Nash if and only if \(x_0\ge x_H(s)\). Below \(x_H\), either the candidate is not branch-feasible or firm \(B\) strictly prefers the no-poaching kink. At equality the candidate remains the unique Nash profile.

The preliminary argument used no nonnegative-price restriction. Therefore the same pure-equilibrium correspondence holds whether uniform net margins range over all real numbers or are restricted to be nonnegative.

\section{Primitive surplus derivations}
\label{app:surplus}

This appendix records the integration logic behind Proposition~\ref{prop:weak-cs}. Normalize \(\tau=1\) and suppress the common gross-utility term, which cancels from differences.

At the uniform source profile,
\[
a_u=1+s/3,
\qquad
b_u=1-s/3,
\qquad
x_u=1/2+s/6.
\]
If \(x_0\ge x_u\), realized uniform consumer surplus is
\begin{align*}
CS^u_{\rm sw}
={}&
\int_0^{x_u}(-a_u-z)\,dz
+
\int_{x_u}^{x_0}[-b_u-(1-z)-s]\,dz\\
&+
\int_{x_0}^{1}[-b_u-(1-z)]\,dz.
\end{align*}
If \(x_0<x_u\), the correct allocation has no switching:
\[
CS^u_{\rm ns}
=
\int_0^{x_0}(-a_u-z)\,dz
+
\int_{x_0}^{1}[-b_u-(1-z)]\,dz.
\]

For weak HBP, substitute the normalized prices in \eqref{eq:weak-hbp-prices}. The within-history cutoffs are
\[
x_A^d=\frac{2x_0+1+s}{6},
\qquad
x_B^d=\frac{2x_0+3-s}{6}.
\]
The primitive integral is
\begin{align*}
CS^d
={}&
\int_0^{x_A^d}(-p_A^d-z)\,dz
+
\int_{x_A^d}^{x_0}[-q_B^d-(1-z)-s]\,dz\\
&+
\int_{x_0}^{x_B^d}(-q_A^d-z-s)\,dz
+
\int_{x_B^d}^{1}[-p_B^d-(1-z)]\,dz,
\end{align*}
where all prices in this display are normalized net margins.

Subtracting the relevant uniform integral and restoring \(\tau\) gives Proposition~\ref{prop:weak-cs}.

\subsection{Sign on the switching branch}

With \(\tau=1\), the numerator of the switching-branch gap is
\[
N_1(x_0,s)
=
-52x_0^2+52x_0+1+36sx_0-34s+s^2.
\]
This is concave in \(x_0\). A switching subinterval inside weak dominance exists only for \(s\le3/5\). Its endpoint values satisfy
\[
N_1(x_u,s)
=
\frac{2(25s^2-72s+63)}9>0,
\]
and
\[
N_1(\bar x,s)
=
\frac{(1+s)(43-45s)}4>0.
\]
The first quadratic has negative discriminant and positive leading coefficient; the second is positive for \(s\le3/5\). Concavity therefore gives \(N_1>0\) throughout the switching part of the weak domain.

\subsection{Sign on the no-switch branch}

The no-switch numerator is
\[
N_0(x_0,s)
=
s^2+12sx_0-14s-8x_0^2+8x_0+5,
\]
also concave in \(x_0\). At the lower endpoint,
\[
N_0(1/2,s)=(1-s)(7-s)>0.
\]
If \(s\le3/5\), the relevant upper endpoint is \(x_u\), where
\[
N_0(x_u,s)
=
\frac{25s^2-72s+63}{9}>0.
\]
If \(s>3/5\), the weak upper endpoint is \(\bar x<x_u\), where
\[
N_0(\bar x,s)
=
\frac{(1-s)(5s+13)}2>0.
\]
Hence the weak profile consumer-surplus gap is positive throughout the source weak-dominance domain.

\subsection{Exact regression points}

Two exact rational checks are retained. At \((s,x_0)=(99/100,501/1000)\), the actual uniform allocation is no-switch and
\[
CS^d-CS^u=\frac{17993}{4500000}>0.
\]
At \((s,x_0)=(1/10,7/10)\), which lies on the one-way-switch branch and inside the corrected pure-overlap domain,
\[
CS^d-CS^u=\frac{221}{720},
\]
while the accepted Eq.~(15) expression evaluates to \(1279/3600\). These checks distinguish the branch-ordering problem from the independent algebraic discrepancy in the accepted display.

\section{Reproducibility and formal scope}
\label{app:reproducibility}

The replication package is designed so that mathematical checks do not require a local copy of the copyrighted source article.

\paragraph{Exact evaluator.}
\texttt{code/uniform\_game\_cleanroom.py} constructs consumer shares directly from history-specific thresholds and clipping. It separately enumerates feasible unilateral best-response candidates and retains exact rational regressions for the Eq.~(12) counterexample, the equality boundary, Eq.~(15)/(16) source fidelity, weak switching/no-switch profit accounting, the strong-domain profit overclaim guard, the \(x_0=1\) empty-segment endpoint, consumer-surplus corrections, welfare identities, and the strong-HBP price-selection family.

\paragraph{Symbolic reconstruction.}
\texttt{code/symbolic\_reconstruction.py} integrates primitive utility and resource costs in SymPy. \texttt{code/downstream\_impact\_symbolic.py} checks the endpoint and factor identities used in the downstream sign proofs.

\paragraph{Generated data.}
\texttt{generated/parameter\_domains.csv}, used in Figure~\ref{fig:domains}, is produced deterministically from the frozen formulas and is compared byte-for-byte with a fresh regeneration in continuous integration.

\paragraph{Formal verification.}
The Lean source in \texttt{formal/Gss2012/Core.lean} checks selected algebraic and inequality statements using Lean 4.19.0 and a pinned mathlib revision. The formal artifact does not encode the complete consumer model or Nash equilibrium definition, so it should be read as an independent certificate for the proof-critical algebra rather than as a machine proof of Proposition~\ref{prop:uniform} from primitive economic definitions.

\paragraph{Source provenance.}
Equation-level source mappings are based on a lawful post-referee accepted manuscript and recorded separately from the executable package. Following an independent submission audit, Eqs.~(13)--(17) were re-transcribed from the source image before the Python and Lean statement maps were recertified. No source PDF is required for the replication suite, and source-facing claims remain version-qualified pending direct comparison with the exact Springer typeset Version of Record.


\paragraph{AI-assisted research workflow.}
Generative AI assistance was used in literature discovery, verification-code drafting, algebraic cross-checking, manuscript structuring and language revision, and research-workflow organization. ChatGPT by OpenAI was used in this project, including GPT-5.6 Sol for workflow/analysis assistance and GPT-6 Astra for an independent pre-submission audit. AI-generated suggestions were not treated as mathematical evidence: the equilibrium claims were re-derived from primitives, exact regressions were run independently of prose claims, source equations were checked against the accepted-manuscript image, and selected proof-critical algebra was separately checked in Lean. The author reviewed the code, claims, citations, and manuscript text and retains full responsibility for the research.

\section*{Declarations}

\paragraph{Funding.}
This research received no external funding.

\paragraph{Competing interests.}
The author declares no competing interests.

\paragraph{Data and code availability.}
No empirical datasets were generated or analysed in this study. The exact-rational, symbolic, formal-verification, and figure-generation code used for the reassessment is included in the accompanying reproducibility repository:
\url{https://github.com/ryotamatsuki/gss-2012-hbpd-welfare-reassessment}.

\paragraph{CRediT authorship contribution statement.}
Ryota Matsuki: Conceptualization; Methodology; Formal analysis; Investigation; Software; Validation; Visualization; Writing -- original draft; Writing -- review and editing; Project administration.

\paragraph{Declaration of generative AI and AI-assisted technologies in the manuscript preparation process.}
During the preparation of this work, the author used ChatGPT (OpenAI), including GPT-5.6 Sol for workflow/analysis assistance and GPT-6 Astra for an independent pre-submission audit, to assist with literature discovery and organization, algebraic and code cross-checking, manuscript structuring and drafting, and language and exposition refinement. AI outputs were not treated as mathematical evidence. The author reviewed and edited the content, rechecked source equations against the accepted-manuscript record, independently verified the mathematical claims and cited sources, and takes full responsibility for the content of the article.


\bibliographystyle{apalike}
\bibliography{manuscript/references}

\end{document}
